% =============================================================================
% TCC2 — SEÇÃO 1: INTRODUÇÃO (REESCRITA COMPLETA)
% =============================================================================
% Este arquivo contém o texto reescrito para o TCC2.
% Instruções: copie cada bloco para o Overleaf, substituindo o correspondente do TCC1.
%
% LEGENDA:
%   [MANTIDO]     = texto idêntico ao TCC1
%   [ATUALIZADO]  = texto com modificações pontuais
%   [NOVO]        = texto completamente novo
%   [REMOVIDO]    = texto do TCC1 que NÃO deve aparecer no TCC2
%   [NOTA]        = observação para você, não incluir no texto final
% =============================================================================


% =============================================================================
% RESUMO  [REESCRITO — substituir integralmente]
% =============================================================================
% [NOTA] O resumo do TCC1 mencionava BERTScore, ROUGE-L e ROUGE-SS.
% Tudo isso foi substituído pelas novas métricas. Também atualizamos os números.

\begin{resumo}
A gestão de avaliações online é um desafio crítico para a reputação de restaurantes. Embora os \textit{Large Language Models} (LLMs) ofereçam uma solução escalável para a geração de respostas, garantir a adequação do tom, da empatia e da consistência da marca em língua portuguesa representa uma lacuna significativa. Este trabalho avalia a viabilidade e a maturidade dos principais LLMs (GPT-4, Gemini e Llama~3) para esta tarefa. Utilizando um corpus expandido de 200 comentários reais de restaurantes, distribuídos em dois períodos temporais (pré e pós-popularização das LLMs), foi conduzida uma análise comparativa de estratégias de engenharia de prompt (\textit{zero-shot} vs.\ \textit{few-shot}) e a simulação de duas personas de marca distintas (sofisticada/formal e acolhedora/informal), totalizando 2.400 respostas geradas em 24 cenários experimentais. A performance dos modelos foi aferida por uma metodologia mista, combinando a avaliação humana qualitativa de 15 avaliadores independentes, com concordância medida por Fleiss' Kappa, Krippendorff's Alpha e ICC, com métricas automáticas multidimensionais: fluência textual via perplexidade (GPT-2 PT-BR), adequação de tom (ToneCal), consistência de marca por similaridade estilométrica (\textit{Writeprints-like}) e empatia percebida (modelo WASSA, em caráter exploratório). Os resultados oferecem um panorama crítico sobre as melhores práticas para a automação de interações com clientes, contribuindo com um \textit{framework} de avaliação multidimensional voltado ao português brasileiro.

\textbf{Palavras-chave:} Engenharia de Prompt, Modelos de Linguagem em Larga Escala (LLMs), Processamento de Linguagem Natural (PLN), Avaliação Multidimensional, Perplexidade, Estilometria, ToneCal.
\end{resumo}


% =============================================================================
% ABSTRACT  [REESCRITO — substituir integralmente]
% =============================================================================

\begin{abstract}
Managing online reviews is a critical challenge for restaurant reputations. Although Large Language Models (LLMs) offer a scalable solution for response generation, ensuring appropriate tone, empathy, and brand consistency in the Portuguese language remains a significant gap. This work evaluates the viability and maturity of leading LLMs (GPT-4, Gemini, and Llama~3) for this task. Using an expanded corpus of 200 real restaurant comments, distributed across two temporal periods (pre- and post-LLM popularization), a comparative analysis was conducted on prompt engineering strategies (zero-shot vs.\ few-shot) and the simulation of two distinct brand personas (sophisticated/formal and welcoming/informal), totaling 2,400 generated responses across 24 experimental scenarios. Model performance was measured using a mixed-methods approach, combining qualitative human evaluation from 15 independent raters --- with agreement assessed via Fleiss' Kappa, Krippendorff's Alpha, and ICC --- with multidimensional automatic metrics: textual fluency via perplexity (GPT-2 PT-BR), tone adequacy (ToneCal), brand consistency through stylometric similarity (Writeprints-like), and perceived empathy (WASSA model, exploratory). The results provide a critical overview of best practices for automating customer interactions, contributing a multidimensional evaluation framework for Brazilian Portuguese.

\textbf{Keywords:} Prompt Engineering, Large Language Models (LLMs), Natural Language Processing (NLP), Multidimensional Evaluation, Perplexity, Stylometry, ToneCal.
\end{abstract}


% =============================================================================
% LISTA DE ABREVIATURAS E SIGLAS  [ATUALIZADO]
% =============================================================================
% [NOTA] Adicionar as novas siglas e remover as que não são mais usadas.
% ADICIONAR:
%   ICC    — Intraclass Correlation Coefficient (Coeficiente de Correlação Intraclasse)
%   TQS    — Tone Quality Score (Pontuação de Qualidade de Tom)
%   WASSA  — Workshop on Computational Approaches to Subjectivity, Sentiment and Social Media Analysis
%
% MANTER:
%   API, BERT, BX, CX, GPT, GPU, ICL, LLM, LSTM, NLP, PLN, RLHF, RNN, UX
%
% REMOVER (se não forem mais citados no texto):
%   BERTScore, LCS, NLTK, OMW, ROUGE, ROUGE-SS


% =============================================================================
% 1 INTRODUÇÃO  [PARÁGRAFOS INICIAIS — MANTIDOS]
% =============================================================================
% [NOTA] Os dois primeiros parágrafos + Figura 1 permanecem IDÊNTICOS ao TCC1.
% São atemporais e bem escritos. Não alterar.

\section{INTRODUÇÃO}

Com o advento das redes sociais, a escolha de um estabelecimento gastronômico passou a ser profundamente impactada por manifestações públicas de comentários, experiências e avaliações publicadas em plataformas como Instagram, TripAdvisor, Google Reviews etc \cite{silva2019}. Assim, consolidando uma dinâmica de influência informacional, onde o usuário sente-se persuadido à pautar-se no conteúdo dessas mídias.

Um exemplo para ajudar a ilustrar esse cenário: em um relato no Instagram de um restaurante, um cliente fez um comentário negativo em relação à demora no atendimento e serviço, avaliando-o como decepcionante: ``Estávamos hospedados em um hotel de alto padrão, mas o restaurante foi decepcionante. O atendimento demorou, o prato infantil veio errado e o meu risoto tinha quase só beterraba, com quase nada de cogumelo. Em um restaurante desse nível, isso é inadmissível. Não voltaria ao restaurante.''. Embora breve, um comentário desse tipo pode influenciar enfaticamente a percepção de futuros clientes. Nesse caso, a resposta do estabelecimento foi: ``Lamentamos muito pela sua experiência, em especial pela troca do prato e pela demora no atendimento. Estamos revisando nossos processos para evitar que situações assim se repitam. Esperamos ter a oportunidade de recebê-lo novamente e oferecer uma experiência condizente com o que valorizamos.''. Visto isso, uma possível resposta de LLMs, sem ajuste, pode soar genérica e antipática: ``Obrigado pelo seu comentário. Sentimos muito pela experiência e valorizamos seu feedback para melhorarmos nossos serviços.''. Essa comparação evidencia o tamanho do desafio e a complexidade na geração de respostas automáticas que soem naturais.

% [NOTA] FIGURA 1: MANTER COMO ESTÁ
% A imagem comment_example.png continua perfeitamente válida.
\begin{figure}[htb]
    \centering
    \includegraphics[width=0.8\textwidth]{assets/comment_example.png}
    \caption{Comentário do cenário exemplificado retirado do Instagram}
    \label{fig:comment_example}
\end{figure}

% =============================================================================
% [MANTIDO] Parágrafo sobre a escala do problema
% =============================================================================

Visto a repercussão que essas avaliações podem causar sobre a reputação de um estabelecimento, o tamanho desse desafio é expressivo. Plataformas como o TripAdvisor acumulam bilhões de avaliações oriundas de diferentes plataformas \cite{tripadvisor2024}. Gerenciar manualmente esse volume, mantendo a qualidade e a agilidade na comunicação, torna-se uma tarefa desafiadora. É nesse cenário que as LLMs apresentam-se como uma ferramenta promissora para a automação dessa tarefa \cite{gamboa2023}. Por um lado, sua capacidade de processar e gerar texto em grande escala oferece agilidade, permitindo que nenhum comentário fique sem resposta. Contudo, em contraponto a essa eficiência, a comunicação humana eficaz demanda sutilezas que vão além de adequação e correção gramatical, exige-se uma percepção de tom para diferenciar uma crítica construtiva de um desabafo, empatia para validar a experiência do cliente e uma compreensão contextual que preserve a identidade do estabelecimento \cite{crolic2023}.

No entanto, o desafio não reside apenas na automação da resposta, mas na capacidade desses sistemas em gerar interações que, de fato, pareçam humanas e inspirem cuidado e confiança. Pesquisas recentes demonstram que a empatia percebida e o tom emocional de agentes de IA impactam diretamente a satisfação do consumidor e a intenção de recompra \cite{rohden2024}, reforçando a criticidade dessa lacuna.

% =============================================================================
% [ATUALIZADO] Parágrafo de síntese — antes "estudaremos", agora reflete TCC2
% =============================================================================
% [NOTA] Mudanças:
%   - "estudaremos" → "este trabalho avalia" (TCC2 é a versão final)
%   - Adicionada menção ao framework multidimensional de avaliação
%   - Menção explícita ao corpus expandido

Em suma, este trabalho avalia a qualidade, viabilidade e maturidade das LLMs disponíveis para a geração de respostas a comentários no setor gastronômico, com foco exclusivo na língua portuguesa, para assim capturar melhor as nuances linguísticas e culturais inerentes ao português \cite{araujo2024}. O estudo compreende um levantamento dos modelos mais destacados, assim como uma análise aprofundada das estratégias de prompt, para adequar variáveis como tom, empatia e consistência de marca. Alinhado à recomendação da literatura por avaliações holísticas e multidimensionais \cite{liang2023,celikyilmaz2021}, a avaliação é conduzida por meio de um \textit{framework} que combina métricas automáticas especializadas --- fluência via perplexidade \cite{pistotti2025}, adequação de tom \cite{danescu2013}, similaridade estilométrica \cite{abbasi2008} e empatia percebida \cite{barriere2023} --- com a avaliação humana de 15 avaliadores independentes, aplicada sobre um corpus expandido de 2.400 respostas geradas. Por fim, busca-se oferecer um panorama crítico e prático sobre os limites e possibilidades dessa tecnologia, contribuindo para um diálogo sobre a automação consciente e eficaz no sensível campo da interação com o cliente.


% =============================================================================
% 1.1 JUSTIFICATIVA  [MANTIDA com ajuste final]
% =============================================================================
% [NOTA] Mudanças:
%   - Último parágrafo: adicionada menção à lacuna em avaliação multidimensional
%   - Resto: idêntico ao TCC1

\subsection{JUSTIFICATIVA}

Estudos demonstram que, no setor gastronômico, as avaliações online impactam diretamente a escolha do consumidor. Essa realidade sujeita os estabelecimentos a uma dura realidade: a necessidade de engajamento constante com o público e o desafio de fazer isso de forma rápida, incessante e, acima de tudo, humana \cite{aureliano2021}. Nesse cenário, as LLMs surgem como uma alternativa escalável e promissora na automação dessas interações, permitindo que empresas gerenciem grandes volumes de dados sem perder a capacidade de resposta operacional. Estudos indicam que fatores como gênero e percepção de lotação podem influenciar significativamente a forma como os consumidores interpretam avaliações e respostas, reforçando a necessidade de respostas contextuais e empáticas \cite{ali2021}. A sua capacidade de compreender perguntas e cenários é o que lhe confere o poder de formular respostas compatíveis, tanto em contexto, quanto em tom, empatia e assertividade. No entanto, essa poderosa ferramenta traz consigo um desafio igualmente relevante: a qualidade da comunicação. Respostas mecânicas, repetitivas ou impessoais podem agravar percepções negativas, comprometendo a imagem do estabelecimento em vez de protegê-la, visto que a falta de calor humano é frequentemente penalizada pelos consumidores \cite{sands2020}.

% [ATUALIZADO] Último parágrafo da justificativa
É justamente nesse ponto que este estudo se faz necessário. Atualmente, observa-se uma lacuna de conhecimento sobre a real qualidade e maturidade das LLMs em português para o setor gastronômico. Embora existam avanços em modelos para o português brasileiro, como o BERTimbau \cite{souza2020}, a aplicação generativa focada em nuances culturais de serviço ainda carece de exploração aprofundada. Além disso, a literatura apresenta uma carência de \textit{frameworks} de avaliação multidimensional que mensurem, de forma integrada e com métricas especializadas, aspectos como fluência, tom, consistência de marca e empatia em respostas geradas automaticamente em português --- uma lacuna corroborada por \citet{liang2023}, que advogam por avaliações holísticas que vão além de métricas agregadas, e por \citet{howcroft2021}, que demonstram que avaliações humanas em PLN são frequentemente subdimensionadas. No campo da empatia computacional, embora \citet{sharma2020} tenham proposto \textit{frameworks} para mensurar empatia em texto e as \textit{shared tasks} da WASSA \cite{barriere2022,barriere2023,giorgi2024} tenham avançado na predição de empatia em reações a textos, a transferência dessas abordagens para o domínio de respostas comerciais em português permanece inexplorada. De forma análoga, a análise estilométrica de consistência de marca, fundamentada em técnicas de \textit{Writeprints} \cite{abbasi2008}, e a avaliação automática de adequação de tom, investigada computacionalmente desde os trabalhos seminais de \citet{danescu2013} e sistematizada por \citet{priya2024}, não foram aplicadas ao contexto de respostas automáticas de restaurantes. Assim, este trabalho propõe-se a ir além da simples automação, conduzindo uma análise comparativa capaz de investigar não apenas a correção linguística, mas também a habilidade dos modelos em modular variáveis como tom, empatia e preservação da identidade da marca. Dessa forma, a pesquisa se justifica por sua relevância prática, ao oferecer diretrizes seguras para o uso consciente dessa tecnologia, contribuindo com um esquema de avaliação multidimensional voltado à automatização de interações em contextos sensíveis.


% =============================================================================
% 1.2 OBJETIVO GERAL  [MANTIDO — idêntico ao TCC1]
% =============================================================================

\subsection{OBJETIVO GERAL}

O objetivo geral deste trabalho é avaliar a viabilidade do uso de LLMs na tarefa de geração automática de respostas para comentários do domínio de restaurantes, com foco exclusivo na língua portuguesa.


% =============================================================================
% 1.3 OBJETIVOS ESPECÍFICOS  [REESCRITO]
% =============================================================================
% [NOTA] Mudanças:
%   - Objetivo 1: menção ao corpus expandido (200 comentários)
%   - Objetivo 4: especifica as 2.400 respostas
%   - Objetivo 5: NOVO — framework de avaliação multidimensional
%   - Objetivo 6: NOVO — avaliação humana com concordância inter-anotador
%   - Removida a menção genérica a "metodologia mista" do TCC1

\subsection{OBJETIVOS ESPECÍFICOS}

Os objetivos específicos deste estudo são:

\begin{itemize}
    \item Coletar, organizar e expandir um corpus de comentários reais de avaliações sobre restaurantes, coletados de plataformas digitais, totalizando 200 comentários distribuídos em dois períodos temporais distintos (pré e pós-popularização das LLMs);

    \item Desenvolver um conjunto de prompts para guiar a geração de respostas das LLMs, explorando variáveis como tom, empatia e alinhamento com a identidade do estabelecimento por meio de duas personas de marca distintas;

    \item Estudar e comparar estratégias de interação com as LLMs (\textit{zero-shot} e \textit{few-shot}), visando maximizar a adequação contextual das respostas e o alinhamento com a identidade do estabelecimento;

    \item Submeter o corpus de comentários e os prompts desenvolvidos aos principais LLMs selecionados para o estudo (GPT-4, Gemini e Llama~3), gerando um conjunto de 2.400 respostas automáticas distribuídas em 24 cenários experimentais;

    \item Desenvolver e aplicar um \textit{framework} de avaliação automática multidimensional, empregando métricas especializadas para fluência (perplexidade via GPT-2 PT-BR) \cite{pistotti2025,serrano2023}, adequação de tom (ToneCal) \cite{danescu2013,balshetwar2019}, consistência de marca (similaridade estilométrica \textit{Writeprints-like}) \cite{abbasi2008} e empatia percebida (modelo WASSA) \cite{barriere2023,giorgi2024};

    \item Conduzir uma avaliação humana qualitativa com 15 avaliadores independentes, utilizando escala Likert \cite{amidei2019}, e aferir a concordância inter-anotador por meio de Fleiss' Kappa, Krippendorff's Alpha e ICC \cite{howcroft2021}.
\end{itemize}


% =============================================================================
% 1.4 ORGANIZAÇÃO DO TRABALHO  [REESCRITO]
% =============================================================================
% [NOTA] Mudanças:
%   - Seção 2: removida ênfase em ROUGE/BERTScore, adicionadas novas métricas
%   - Seção 4: removida "transição ROUGE-SS→BERTScore", atualizada com novas métricas
%   - Seção 5: atualizada para refletir análise multidimensional
%   - Seção 6: "Conclusões" (não mais "Conclusões Parciais", pois TCC2 é versão final)

\subsection{ORGANIZAÇÃO DO TRABALHO}

O presente trabalho está estruturado em seis seções principais, organizadas para guiar o leitor desde a fundamentação teórica até a análise dos resultados experimentais.

A Seção~1 (Introdução) contextualiza o problema da gestão de avaliações online no setor gastronômico, apresenta a justificativa para o uso de LLMs na automação de respostas e define os objetivos gerais e específicos da pesquisa.

A Seção~2 (Fundamentação Teórica) estabelece as bases conceituais necessárias para o estudo, abordando o estado da arte das LLMs, os princípios da Engenharia de Prompts (incluindo estratégias \textit{Zero-Shot} e \textit{Few-Shot}) e as métricas de avaliação de texto gerado, com ênfase nas abordagens de fluência, adequação de tom, consistência estilométrica e empatia percebida.

A Seção~3 (Trabalhos Relacionados) revisa a literatura acadêmica pertinente, destacando estudos anteriores sobre geração automática de respostas, aplicações em atendimento ao cliente e as lacunas existentes no processamento de linguagem natural para o português brasileiro.

A Seção~4 (Metodologia) detalha a arquitetura experimental adotada, descrevendo a coleta e expansão do corpus de comentários, o desenvolvimento dos prompts e das personas, a justificativa metodológica para a seleção dos modelos (GPT-4, Gemini e Llama~3), o \textit{framework} de avaliação automática multidimensional e o protocolo de avaliação humana com 15 avaliadores e medidas de concordância inter-anotador.

A Seção~5 (Experimentos e Resultados) apresenta e discute os dados obtidos a partir dos 24 cenários experimentais. Esta seção analisa o desempenho dos modelos sob a ótica de quatro dimensões automáticas --- fluência, adequação de tom, consistência de marca e empatia percebida --- complementadas pela avaliação humana qualitativa.

Por fim, a Seção~6 (Conclusões) sintetiza as descobertas do estudo, validando as hipóteses iniciais sobre a viabilidade técnica e qualitativa da automação, e aponta as contribuições, limitações e direções para trabalhos futuros.


% =============================================================================
% 1.5 CRONOGRAMA DAS ATIVIDADES  [ATUALIZADO]
% =============================================================================
% [NOTA] Atualizar a tabela do cronograma para refletir as atividades do TCC2.
% O TCC1 já foi feito. Agora o foco é na segunda etapa.
% Ajuste os meses conforme o calendário real do seu TCC2.
% Sugestão de atividades atualizadas:

\subsection{CRONOGRAMA DAS ATIVIDADES}

A Tabela~\ref{tab:cronograma} apresenta o cronograma das atividades realizadas ao longo das duas etapas deste trabalho.

% [NOTA PARA O OVERLEAF] Recrie a tabela com as seguintes atividades:
%
% Atividade                              | Meses (adaptar ao seu calendário real)
% ----------------------------------------|----------------------------------------
% Coleta de comentários (TCC1)            | mês 2
% Engenharia de prompts (TCC1)            | mês 3
% Submissão às LLMs (TCC1)               | mês 4-5
% Escrita e defesa TCC1                   | mês 6-7
% Expansão do corpus                      | mês 8
% Geração das 2.400 respostas             | mês 8-9
% Avaliação humana (15 avaliadores)       | mês 9-10
% Implementação métricas automáticas      | mês 9-11
% Análise e consolidação dos resultados   | mês 11-12
% Escrita TCC2                            | mês 8-12
% Defesa TCC2                             | mês 12


% =============================================================================
% FIGURAS DA INTRODUÇÃO — ANÁLISE
% =============================================================================
%
% FIGURA 1 (comment_example.png):
%   ✅ MANTER — screenshot do Instagram mostrando comentário real + resposta humana.
%   Continua sendo um exemplo perfeito e atemporal do problema.
%   Nenhuma alteração necessária.
%
% [NOTA] As Figuras 2, 3 e 4 pertencem à Seção 2 (Fundamentação Teórica):
%
%   FIGURA 2 (comentário_exemplo_intro.png):
%     ✅ MANTER — exemplo do aniversário com resposta empática.
%     Ilustra bem a diferença entre resposta LLM genérica e resposta humana acolhedora.
%
%   FIGURA 3 (comentário_resposta_genérica.png / resposta_genérica.png):
%     ✅ MANTER — exemplo de resposta genérica e dismissiva ("Olá, tudo bem?
%     Enviamos uma mensagem para você!"). Excelente exemplo de falha real.
%
%   FIGURA 4 (comentário_resposta_inconsistente.png):
%     ✅ MANTER — exemplo de inconsistência de marca ("Opa, tudo certo?
%     Uma pena que você não curtiu a vibe" em restaurante requintado).
%     Ilustra perfeitamente o problema de consistência de marca.
%
% CONCLUSÃO SOBRE FIGURAS DA INTRODUÇÃO:
%   → Nenhuma figura precisa ser alterada ou substituída.
%   → Todas continuam relevantes e bem contextualizadas.
%   → Na Seção 5 (Resultados), novas figuras serão necessárias
%     (histogramas de fluência, gráficos de ToneCal, etc. — já existem nos results/).
